% !TeX root = ../main.tex

\chapter{绪论}

\section{研究背景及意义}

识别网络流量类型是开展网络空间安全研究与管理的一项基础且核心的能力。流量分类旨在将复杂多变的网络数据流准确归类至其对应的特定业务类别中，这对于实现服务质量优化、差异化计费、网络资源规划、恶意代码识别以及入侵检测等关键应用具有不可替代的支撑作用。

随着用户隐私保护意识的显著提升，加密通信协议已逐渐成为互联网流量的主导载体。根据 Firefox 的统计数据，美国互联网用户利用 SSL/TLS 安全协议进行通信的比例已从 2017 年的 55\% 激增至 2026 年的 95\%。与此同时，虚拟专用网络（VPN）和洋葱路由（Tor）等匿名化技术的广泛普及，进一步强化了网络通信的隐蔽性。在此趋势下，如何在不破坏端到端加密原则的前提下实现精准的流量分类，已成为保障网络安全供给与提升用户体验质量的关键技术。

\begin{figure}[ht]
  \centering
  \includegraphics[width=1\textwidth]{firefox.png}
  \caption{加密协议使用情况}
  \label{fig:firefox}
\end{figure}

早期的流量分类方案主要依赖于端口识别技术和深度包检测技术（Deep Packet Inspection，DPI）\cite{nelms2013execscent,su201660gbps,melo2014performance}。基于端口的方法通过匹配传输层协议端口与应用服务的对应关系进行识别，但随着动态端口协商及端口转发技术的广泛应用，该方法已逐渐失效。DPI 技术通过解析 IP 数据包的完整载荷以提取特征，然而在加密环境下，数据包载荷呈现为高熵的乱序字节序列，DPI 因无法穿透加密层获取明文特征，同样面临严峻的失效困境。

人工智能技术的进步为加密流量分类注入了新的生命力，主要围绕着传统机器学习和深度学习算法的应用。基于机器学习的方法通常基于人工预设的统计特征构建流量特征，如数据包路径签名、数据包到达时间和大小以及累积数据包大小等，并且人工选择性能最佳的分类器对其进行分类。然而，机器学习方法的性能高度依赖于特征工程的质量，即为每个特定的加密流量协议或工具选择最优特征，而选择优质的特征表示要求研究者具备深厚的领域专家知识。此外，由于人工提取的特征往往针对特定协议设计，其泛化能力较差，难以适应现网中层出不穷的新型加密工具与协议。在真实的网络场景中，为不同工具或协议生成的各类加密流量训练和实施不同的基于机器学习的加密流量分类方法既耗时又难以实现。 

作为机器学习的一个分支，深度学习在加密流量分类领域展现出巨大的潜力。深度学习技术的核心优势在于：模型能够直接从原始流量数据中自动执行多维度的表征学习，消除了对复杂特征工程的依赖。且相比于依赖完整流统计信息的机器学习方法，深度学习模型能够捕捉报文交互初期的细微特征，从而在流量会话早期即完成识别。此外，深度学习采用端到端的学习范式，通过深层神经网络架构，直接从原始输入数据学习以产生期望的输出，绕过了将问题分解为特征选择和分类等子问题的传统做法，能够更全面地挖掘流量数据中隐含的复杂关联模式，这一特性使深度学习算法能够全面地处理流量分类任务相关的复杂非线性问题。

尽管基于深度学习的方法已取得显著进展，但通过对现有主流算法的调研发现，仍存在许多有待解决的问题。首先多数算法倾向于单一维度分析，仅关注由包长度序列构成的流级别特征，或仅提取数据包内容构造包级别特征，未能充分挖掘跨维度特征间的协同表征能力，使得流量特征表示不充分。其次，部分研究在提取包级别特征时，未对报文首部与有效载荷的内容进行功能解耦，将其作为统一整体处理，忽略了两者在协议语义与信息载荷上的显著差异。最后，部分模型对有效载荷数据的深度依赖不仅增加了计算负载，也可能在处理敏感数据时违反用户隐私保护原则。综上所述，虽然基于深度学习的加密流量分类已成为主流研究方向，但如何高效提取具备泛化能力的特征，并充分利用多维度流量信息，仍是当前面临的核心挑战。

\section{国内外研究现状}
随着全球互联网信息化进程的加速以及数字经济的蓬勃发展，网络流量的爆发式增长使得流量分类技术成为网络管理、服务质量保障及网络安全监管领域的核心课题。经过数十年的演进，流量分类技术已从早期的简单匹配发展为如今基于人工智能的高度自动化识别。本文将从基于载荷检测、基于传统机器学习以及基于深度学习三个维度，对国内外现有的研究成果进行介绍。

\subsection{基于载荷检测的方法}

基于载荷检测的流量分类方法，也称为深度包检测（DPI），该技术的核心逻辑在于对数据包的应用层载荷进行拆解，并将其内容与预定义的协议特征签名进行模式匹配\cite{khalife2014multilevel}。常用的高性能DPI库包括libprotoident、OpenDPI和nDPI。其中，nDPI 通过内置针对 HTTP、FTP 等主流协议的专用解码器，实现流量识别\cite{xie2020sam}。除有效载荷内容外，部分 DPI 方案也会结合报文首部字段进行辅助判断，其中最常用的字段是端口号。例如，目的地为 80 端口的数据包通常会被优先分流至 HTTP 解析模块\cite{constantinou2006identifying,erman2007identifying}。然而，这种方法在端口动态分配和通用通信协议端口的情况下会失效。

DPI 技术目前在实际部署中面临严峻挑战。首先是计算开销问题，由于涉及海量字符串的逐一比对，其处理性能往往难以匹配高速骨干网的需求。尽管Khandait等人提出了一种基于DPI的网络流量分类方法，使用启发式方法，通过单次扫描流载荷即可对网络流进行分类，实现了次线性搜索复杂度，但在处理大规模并发流量时依然压力显著\cite{khandait2020efficient}。其次是局限性与合规性问题，随着 SSL/TLS 加密协议的普及以及各类私有协议的兴起，载荷内容变得不再透明，导致特征匹配机制失效 。此外，深度的有效载荷提取往往触及用户隐私保护的法律红线，限制了其在特定场景下的应用\cite{dainotti2012issues}。

\subsection{基于传统机器学习的方法}
随着机器学习技术的出现，研究者开始转向基于统计特征的机器学习方法。此类方法的核心在于特征工程，即通过人工经验提取流量的统计属性，并利用监督学习算法构建分类器，其分类准确性的最关键因素变成了如何构造足够有效的特征并选择性能最优的分类器。

Taylor 等人开发的 AppScanner 利用随机森林算法对流数据的统计特征进行建模，实现了对移动端流量的有效区分\cite{taylor2016appscanner}。Liu 等人提出了一种多属性马尔可夫概率模型，该模型整合了消息类型序列和数据包长度，用于对流量应用进行分类\cite{liu2018mampf}。Ede 等人提出了一种名为 FlowPrint 的移动应用识别方法，通过半监督学习手段，自动查找网络流量中与目标相关的特征之间的时间相关性，并利用这些相关性生成应用指纹\cite{van2020flowprint}。Taylor 等人探讨了被动窃听者如何通过分析智能手机应用传输的网络流量来对其进行识别。尽管数据包的有效载荷被 SSL/TLS 加密所隐藏，利用基于数据包大小和方向等侧信道数据仍可以识别智能手机应用。此外还进一步研究了应用程序指纹如何随时间、跨设备以及跨不同版本的应用程序而变化\cite{taylor2017robust}。Shen 等人提出了一种名为 FineWP 的网页指纹识别方法，利用双向客户端-服务器交互中数据包的长度信息提取出作为网页指纹识别的独特特征，并将其与上行主导阶段的数据包位置相结合，随后输入到随机森林、决策树和 k 近邻等机器学习算法中以创建细粒度的网页指纹\cite{shen2020fine}。Saber 等人提出了一种基于主成分分析的方法用于流量分类，该方法仅关注基于时间的数据流流特征，并利用支持向量机分类器在进行高效流量分类之前选择最佳特征子集\cite{saber2018encrypted}。Anderson 等人将流元数据、包长度分布、时间分布、字节分布以及未加密的TLS头部信息作为联合特征，利用逻辑回归算法来识别恶意软件加密流量\cite{anderson2018deciphering}。Liu等人设计了包级统计特征，利用发送和接收字节的最大值、最小值和平均值作为原始数据进行特征提取，并提出了一种基于复合特征的半监督方法用于加密流量分类\cite{liu2012semi}。Zhang等人提出了一种新的鲁棒流量分类（RTC）方案，该方案使用基于词袋（BoF）的流量分类器对数据流进行处理，并通过实验证明其性能显著优于最常见的机器学习方法\cite{zhang2014robust}。Fu等人将加密流量流以分层方式分割成会话，并提取包长度和时间延迟序列来构建隐马尔可夫模型，该模型能够识别服务使用情况和终端用户的应用程序内行为\cite{fu2016service}。Feghhi等人提出仅使用上行链路的数据包时序信息来测量 VPN 流量的方法，该方法无需了解网页请求的起始/结束时间，且不受数据包填充防御机制的影响\cite{feghhi2016web}。Panchenko等人首次提出利用累计数据包大小特征生成Tor网站指纹，并被后续研究者广泛采用\cite{panchenko2016website}。

以上基于机器学习的方法都存在一个共性，其粒度都在流或会话级别，因为要获得通信模式并提取统计特征，需要观察整个流或会话，等待完整流结束才能进行决策，较高的时延使其无法满足实时分类的需求。同时，基于机器学习的方法都建立在特征工程的基础上，分类效果高度依赖专家的先验知识，且针对特定协议选取的特征往往缺乏通用性，在面对协议更新或流量分布偏移时，模型的鲁棒性较差。因此，基于机器学习的方法无法快速统一地选择的特征来分类加密流量，不适合流量类型和来源多变的真实网络场景。

\subsection{基于深度学习的方法}

受深度学习在图像处理和自然语言处理领域卓越性能的推动，研究人员将深度学习应用于网络流量分类。基于深度学习的方法不同于传统的基于机器学习的方法，其能够直接从原始字节流中提取高维非线性特征，有效克服了传统方法对专家经验的过度依赖。

早期的深度学习尝试多聚焦于将流量进行图像化，Wang等人提出了一种用于加密流量分类的端到端方法，该方法将原始流量的字节数据作为输入，并将这些数据转换为灰度图像，再用一维CNN进行处理分类\cite{wang2017end}。Shapira 等人提出了一种用于加密互联网流量分类和应用识别的新方法FlowPic，该方法将基本的流数据转换为直观的图像，然后使用卷积神经网络来识别流类别和正在使用的应用程序\cite{shapira2021flowpic}。Liu 等人提出了一个通过深度神经网络进行加密流量分类的框架FS-Net，将数据包长度序列作为输入，采用双向门控循环单元(Bi-GRU)进行特征编码，并在自动编码器中结合了重构机制以保证学习到的特征的有效性\cite{liu2019fs}。Liu 等人提出了一个名为 BGRUA 的模型，用于识别在 HTTPS 连接上运行的 Web 服务。BGRUA 同样采用双向门控循环单元，从会话内的字节序列中捕获前向和后向特征。此外，还根据特征对分类任务的贡献程度为其分配不同权重\cite{liu2020attention}。Lin 等人提出了一种新的流量表示模型，称为加密流量双向 Transformer 编码器表示(ET-BERT)。该模型利用大规模未标记数据进行预训练，并根据两项预训练任务提取具有强力上下文联系的数据报级表示。第一项任务是通过预测被遮蔽的令牌来捕获流量字节间的上下文关系。第二项任务则是通过预测同源突发流来预测正确的传输顺序\cite{lin2022bert}。Lotfollahi 等人提出了一种基于深度学习的加密流量分类方法Deep Packet。该方法利用堆叠自动编码器和卷积神经网络进行应用识别任务\cite{lotfollahi2020deep}。Aceto 等人提出了一种名为 MIMETIC 的多模态深度学习方法。他们分别利用 CNN 和 GRU 来学习两种不同的模态的特征，即有效载荷的前 576 个字节和信息性协议字段，并将其用于移动加密流量分类\cite{aceto2019mimetic}。Wang 等人也提出了一个类似的模型 App-Net，该模型使用 LSTM 和 CNN 分别学习数据包长度序列和初始数据包的有效载荷字节数据，分别提取特征\cite{wang2020app}。图神经网络 (GNN) 在处理非结构化数据方面具有强大潜力，并在许多领域展现出良好的性能。Shen 等人提出了 GraphDApp，该模型能够使用 GNN 识别去中心化应用 (DApps) 的加密流量。他们基于数据包长度和数据包方向构建了加密流的流量交互图 (Traffic Interaction Graph，TIG)，并将 DApp 识别问题转化为图分类问题\cite{shen2021accurate}。Han等人提出了双嵌入图神经网络（DE-GNN），分别用 one-hot 编码处理包头和载荷数据，并设计PackertCNN提取包级别特征，然后构建流量交互图（TIG）并使用图神经网络学习流级别特征，最后用自适应特征融合将二者融合作为最终的模型训练特征用于训练分类\cite{han2024gnn}。Yang等人提出了一个利用双模态特征的DM-HNN模型，该模型同时使用包长度序列作为流级别特征和数据包数据作为包级别特征，并通过门控激活机制自适应调整特征权重将其融合作为最终用于分类的特征进行模型训练\cite{yang2023network}。

现有深度学习方案虽已取得显著成效，但仍存在一定的局限性。首先，多数研究仍局限于单一模态，如仅关注包长度序列或仅关注有效载荷，未能充分挖掘流级特征与包级特征之间的协同效应。而现有的多模态融合尝试多为简单的向量拼接，没有考虑到特征之间的相互关系\cite{wang2024combine}。此外，基于包长度序列的方法往往忽略了复杂网络环境如抖动、丢包等对时序特征的干扰，从而导致特征表达不充分。而使用包级别特征的方法依赖有效载荷数据的使用，这在增强分类精度的同时，不可避免地增加了隐私泄露风险。

\section{研究内容}
针对上述问题，本文提出了两项研究，首先针对加密流量特征表达不充分、鲁棒性缺失等问题，本文研究并提出了一种基于自监督学习框架 BYOL 的流级别特征分类方法，该方法立足于流级别特征数据，通过构造抗干扰的特征提取模型，捕捉在网络波动影响下仍然不变的流级别鲁棒性特征，完成在复杂网络环境下也能顺利分类的任务。接着在获取高性能流级别特征的基础上，本文针对隐私保护不足、特征维度单一以及单一模态特征存在的类别混淆问题，进一步研究并提出了一种基于融合特征的加密流量分类方法，通过提取数据包内容构造包级别特征，并将两个模态的特征进行融合，从而构建具备高分类性能与强鲁棒性的表示特征，进一步增强了加密流量分类模型的有效性。通过上述两个阶段的研究，本文构建了从底层特征自动化表征到多维度特征深度融合的模型架构。具体研究框架如图 \ref{fig:content} 所示，主要研究内容涵盖以下三个方面：

\begin{figure}[ht]
  \centering
  \adjincludegraphics[width=1\textwidth, trim=2 2 2 2, clip]{content.pdf}
  \caption{研究内容结构图}
  \label{fig:content}
\end{figure}

（1）面向网络环境波动的流级别鲁棒特征提取。
针对网络拥塞、路径抖动及分片机制导致的流量时序特征失真问题，本文研究如何构建具备抗干扰能力的流级别特征表征模型。针对不同网络状态下的统计特性与时序规律提出不同的增强技术，并利用自监督对比学习机制捕捉流量中不随环境波动的不变性特征，缓解网络波动因素对分类模型泛效化能力的影响，提升模型在真实复杂场景下的分类可靠性。

（2）兼顾表征能力与隐私保护的包级别特征提取。
针对传统方法在处理包级别特征时，因过度依赖应用层有效载荷而带来的隐私安全与合规性问题，本文探索了一种仅基于报文首部信息的特征提取策略。研究重点在于如何在完全规避敏感负载数据的前提下，通过对IP报文首部中协议号、标志位及窗口大小等结构化的关键字信息进行提取，并采取n-gram嵌入处理，捕捉数据报头的字节结构信息、文本顺序信息以及其与特定业务类型之间的隐式关联，从而在确保用户隐私合规性的同时，维持高水平的特征表示。

（3）基于自适应注意力机制的双模态特征融合机制。
针对单一模态特征表达能力受限以及传统融合策略，如简单拼接、静态加权难以有效处理特征间非线性关系的问题，本文研究并构建了一种基于注意力机制动态分配权重的双模态特征融合框架。通过引入注意力机制，模型能够根据流级别与包级别特征在不同业务场景下的贡献度差异，自适应地调整各模态特征在融合过程中的融合比例，从而充分利用双模态的互补信息，实现更深层次的特征语义对齐。

\section{章节安排}
本文主要包括五章，各章节主要内容的总结如下：

第一章为绪论。首先详细阐述了在网络空间安全日益严峻及加密流量普及的背景下，开展加密流量分类研究的重要现实意义与应用价值。接着，对国内外在该领域的研究现状进行了系统综述，对比分析了基于端口检测、深度包检测、传统机器学习以及现有深度学习方法的优缺点，指出当前研究在流级与包级特征融合、特征表示不足及隐私保护等方面存在的局限。在此基础上明确了本文的研究目标，并概括了论文在提升特征鲁棒性、强化隐私保护及构建动态融合框架等方面的三个主要研究内容，最后给出了全文的组织架构。

第二章为相关理论概述。本章为后续研究构建了必要的理论支撑体系，首先回顾了网络流量分类的基础理论，深入剖析了端口检测、DPI及传统机器学习方法的内在机制与性能瓶颈。随后，重点探讨了深度学习的核心架构，详细介绍了卷积神经网络的层级结构、各类非线性激活函数与损失函数的作用机理，以及残差网络解决退化问题的优势。此外，本章还详细阐述了注意力机制在特征增强中的应用，以及自监督学习特别是对比学习在无标签数据表征学习中的前沿理论，为后文模型的设计奠定了坚实的数学与算法基础。

第三章为基于BYOL模型的流级特征分类方法。针对网络环境波动，如拥塞、抖动、分片等导致流量包长度序列失真的难题，本章提出了一种基于自监督对比学习的流级别鲁棒特征提取模型。通过引入BYOL框架，并设计了专门针对流量时序数据的增强函数，使得模型能够在无需负样本的情况下，通过在线网络与目标网络的相互映射学习，捕捉流量中不随环境变化的一致性表示。本章详细介绍了该模型的网络架构、编码器、投影器与预测器设计，并通过多组对比实验验证了该方法在提升流级别特征稳定性及分类精度方面的卓越表现。

第四章为基于双模态融合特征的加密流量分类研究。在流级特征基础上，为解决单一维度特征表达能力不足导致的类别混淆问题，本章进一步探索了隐私保护导向的包级别特征提取策略，仅利用报文首部结构化信息作为原始输入，并通过n-gram嵌入捕捉字节之间的结构信息与顺序信息。此外还构建了一个基于自适应注意力机制的双模态融合框架，利用跨模态注意力机制动态量化流级别时序特征与包级别结构特征在不同业务场景下的贡献度，实现了特征间的深度语义对齐与非线性映射融合。实验结果表明，该融合模型在处理高度混淆的加密流量时，其准确率与可靠性均优于单一模态或简单拼接的分类方法。

第五章为总结与展望。本章首先对全文的研究工作进行了系统性总结，回顾了从自监督表征学习到多模态特征融合的完整加密流量分类体系的构建过程，并再次强调了本文在提升特征鲁棒性与隐私合规性方面取得的研究成果。随后，客观分析了当前模型在应对超大规模网络环境及异构网络自适应能力方面的局限性。最后，针对网络环境的动态化与智能化趋势，提出了未来研究的几个重要方向，包括加强跨域迁移学习能力、优化模型推理效率以及探索更精细化的流量行为建模技术。